% =====================================================================
% 第 3 章  数据与度量
% ---------------------------------------------------------------------
% 编译说明与第 5 章一致：ctexart 下用 \section，ctexbook/report 下改为 \chapter。
% 依赖宏包：booktabs、threeparttable、amsmath
% =====================================================================

\section{数据与度量}

本章说明三类数据的来源与构造：报告文本及其情绪指标、行业行情与市场组合、
事件日历与收益窗口。三者的对齐方式决定了后续识别是否成立，因此对时间戳的处理
说明得较为详细。

\subsection{报告文本与情绪度量}

\paragraph{语料来源与预处理}
研究对象为人民银行按季发布的《中国货币政策执行报告》，样本覆盖 2001 年第一季度至
2026 年第二季度共 102 期，全部取自人民银行官网公开页面。每期报告以 PDF 形式存档，
平均 54 页（最短 11 页、最长 73 页）。文本抽取采用 pypdf，随后做三类清洗：去除页眉
页脚与单独成行的页码；合并被换行截断的句子（2001--2003 年的 PDF 在中文字符之间插入了
空格，需先删除）；剔除附表与专栏目录之后的统计附录。

清洗后单期正文平均 25481 个汉字，中位数 26551，最短 9762（2003Q2），最长 35208。
早期报告篇幅明显短于近年，这一点在解释分段长度时需要留意。

\paragraph{情感词典}
情绪度量采用姜富伟等（2020）发布的中文金融情感词典。该词典针对中文财经文本构建，
经人工筛选与词向量扩充，包含积极词 3338 个、消极词 5890 个，使用需引用
Jiang et al.（2019）与姜富伟等（2020）。考虑到本文采用子串匹配识别情感词，
单字词容易在无关词语中被误命中（例如“好”出现在“好处”“良好”之外的搭配中），
因此剔除 4 个单字积极词与 17 个单字消极词；去重后实际使用积极词 3190 个、
消极词 5604 个。

\paragraph{情感单元与计分}
计分以句子为基本单元。把段落按句号、问号、感叹号与分号切分，在每个句子内部自左向右
扫描，优先匹配较长的词（最多四字），命中后跳过该词长度以避免重复计数。若命中词之前
三个字符内出现否定词（“不”“未”“无”“没”“非”等），则翻转其极性。
单期情绪定义为
\begin{equation}
\text{tone}=\frac{N_{\text{pos}}-N_{\text{neg}}}{N_{\text{pos}}+N_{\text{neg}}},
\end{equation}
其中 $N_{\text{pos}}$、$N_{\text{neg}}$ 分别为翻转极性后的积极词与消极词计数。
该比值落在 $[-1,1]$，取值越高表示措辞越偏宽松。

\paragraph{政策指引段与宏观回顾段的切分}
既有研究指出，报告中对市场起作用的主要是面向未来的政策指引，而非对已公布宏观数据的
回顾（姜富伟等，2021）。为区分二者，本文按报告的章节结构切分：以“货币政策趋势”
（2005 年以后）、“预测与展望”或“预测和展望”（2002Q4--2004Q1）、“下一阶段……”
（2001--2002 年部分报告）作为锚点，锚点之后的正文记为政策指引段（记为 P 段），
之前的部分记为宏观回顾段（记为 R 段）。锚点取全文最后一次出现的位置，以避免命中目录；
若 P 段之后出现“附录”“附表”等标题，则在其处截断。

分段结果见表~\ref{tab:corpus}。P 段平均 3216 字，约占全文的 12.7\%，与报告的章节
安排一致。R 段与 P 段的情绪相关系数为 0.375，说明二者相关但并不同质，分段是必要的。

\begin{table}[htbp]
\centering
\small
\begin{threeparttable}
\caption{报告语料与分段统计}
\label{tab:corpus}
\begin{tabular}{lcccccc}
\toprule
 & 均值 & 标准差 & 最小值 & 中位数 & 最大值 & 观测数 \\
\midrule
全文字数            & 25481 & 4833 & 9762  & 26551 & 35208 & 102 \\
宏观回顾段字数（R） & 22264 & 4436 & 1045  & 22858 & 29765 & 102 \\
政策指引段字数（P） & 3216  & 1737 & 0     & 3122  & 15300 & 102 \\
P 段占全文比重      & 0.127 & 0.089 & 0.000 & 0.120 & 0.936 & 102 \\
P 段情绪（tone\_P） & 0.632 & 0.087 & 0.392 & 0.634 & 0.853 & 99  \\
R 段情绪（tone\_R） & 0.524 & 0.085 & 0.275 & 0.530 & 0.739 & 102 \\
\bottomrule
\end{tabular}
\begin{tablenotes}[flushleft]
\footnotesize
\item 字数指清洗后的汉字数。tone\_P 与 tone\_R 为未标准化的原始比值。
R 段与 P 段情绪的相关系数为 0.375。
\end{tablenotes}
\end{threeparttable}
\end{table}

\paragraph{三项质量说明}
第一，词典覆盖率。政策指引段每千字平均命中 103.8 个情感词，宏观回顾段为 65.0 个，
说明词典对政策表述的识别能力并不弱于对经济描述。

第二，分段失败的处理。102 期中有 3 期（2003Q2、2003Q3、2004Q1）因报告结构差异未能
定位到可靠锚点，未生成 P 段情绪，不纳入回归。另有 1 期（2003Q1）分段后 P 段占比达
93.6\%，明显异常，复核系文本排版造成锚点错位。该期情绪值处于样本中部，剔除后主要
结论不变（见第 5 章稳健性检验），故在主回归中保留并予以说明。

第三，与既有研究的关系。本文的分段依据报告章节，而姜富伟等（2021）按句子内容分类，
两者并不完全等价。这是本文与原文在度量上的主要差异，也是本文结果在方向上与之不完全
一致的可能原因之一，第 6 章将对此作进一步讨论。

\subsection{行业收益与市场组合}

\paragraph{行情数据}
行业行情来自 akshare 提供的申万一级行业指数日线，共 37 个代码，覆盖
1999 年 12 月 30 日至 2026 年 11 月 11 日，合计 6498 个交易日。数据字段包括开高低收
与成交量额。

需要说明的是，37 个代码并非同一套分类。申万于 2014 年与 2021 年两次修订行业分类，
新代码的历史被重置，因此各行业的可用起点分为三档（表~\ref{tab:industry}）。
档 A 的 22 个行业中，5 个（801060、801090、801100、801190、801220）在 2014 年 2 月
停止更新，1 个（801020）在 2021 年 12 月终止，其余 16 个贯穿全样本。

\begin{table}[htbp]
\centering
\small
\begin{threeparttable}
\caption{申万一级行业的数据起点}
\label{tab:industry}
\begin{tabular}{lcl}
\toprule
起点 & 行业数 & 说明 \\
\midrule
1999-12-30 & 22 & 其中 5 个于 2014 年 2 月停止，1 个于 2021 年 12 月终止 \\
2014-02-21 & 12 & 申万 2014 版分类新增或更名的行业 \\
2021-12-13 & 3  & 石油石化、环保、美容护理 \\
\midrule
连续覆盖全样本 & 16 & 本文构造市场组合所使用者 \\
2014 年后覆盖 & 28 & 稳健性检验所使用者 \\
\bottomrule
\end{tabular}
\begin{tablenotes}[flushleft]
\footnotesize
\item 行业分类以申万一级为准，代码与名称的对应关系见附录。
\end{tablenotes}
\end{threeparttable}
\end{table}

\paragraph{市场组合的构造}
本文以 16 个连续行业的等权组合作为市场代理：对每个交易日，取当日有行情的行业日收益
的平均值，再累乘为指数。选择等权而非市值加权，主要出于两点考虑：一是收盘价数据中
不含流通市值，无法可靠加权；二是等权组合与后文的行业面板口径一致，便于对照。

这一处理有代价。等权组合偏向中小市值行业，对大盘蓝筹的代表性弱于沪深 300。
本文未能取得覆盖 2001 年以来的沪深 300 日线，因而以行业组合替代，并在 5.5 节用
28 个行业的口径检验结论对行业集合的敏感性。读者在解读效应量级时应考虑这一差异。

\paragraph{收益与波动的定义}
设市场指数在交易日 $t$ 的取值为 $I_t$。发布后一日收益为 $R^{(1)}=I_{t_1}/I_{t_0}-1$，
五日收益为 $R^{(5)}=I_{t_5}/I_{t_0}-1$，其中 $t_0$、$t_1$、$t_5$ 的定义见 3.3 节。
累计异常收益取 $I_{t_1}/I_{t_{-1}}-1$，即包含发布前一个交易日在内的三日窗口。

波动指标有两个。一是事件日绝对收益 $|R^{(0)}|=|I_{t_0}/I_{t_{-1}}-1|$，计算简便但
无法区分单边移动与日内震荡。二是日内振幅，定义为各行业当日最高价与最低价之差除以前
收盘价，再对 16 个行业取平均。相比绝对收益，振幅不受跳空缺口影响，是更干净的波动代理。
此外，本文以发布前 60 个交易日市场收益的标准差衡量事前波动，用于控制波动聚集。

\subsection{事件窗口与时间对齐}

报告的发布日与报告期并不一致。102 期报告全部在报告期结束后的下一个季度发布，
例如 2026 年第一季度报告于 2026 年 5 月 11 日发布。若把报告期收益与当期情绪放在
同一行，等于用 5 月已知的信息解释 1 至 3 月的行情，构成前视偏差。因此本文以发布日
为事件起点，全部收益与波动指标都在发布之后度量。

事件日 $t_0$ 取发布日当日；若发布日不是交易日，则顺延至其后第一个交易日。
102 期中有 6 期需要顺延：2010Q4、2013Q4、2015Q4、2019Q3、2020Q1、2022Q1。
其中 2015Q4 最为特殊，报告于 2016 年 2 月 6 日（周六）发布，其后是春节休市，
第一个交易日为 2016 年 2 月 15 日，距发布已 9 天。报告发布与市场开盘之间的信息
积累会集中释放，该期的事件窗收益不宜与其他期直接比较，本文在稳健性检验中单独标注。

以 $t_0$ 为起点，$t_1$、$t_5$、$t_{-1}$ 分别为其后第 1 个、第 5 个交易日与之前第
1 个交易日。样本中有 10 期的事件窗跨越 9 个以上自然日（春节或国庆），$t_0$ 至 $t_5$
的自然日跨度中位数为 7 天。所有窗口均按交易日推进，不使用自然日，以免在长假期间
产生空窗口。

\subsection{变量定义与样本}

表~\ref{tab:var} 汇总主要变量。情绪变量均作标准化处理，均值为 0、标准差为 1，
因此回归系数可直接读作情绪每变动一个标准差的效应。制度断点变量 $\text{post}_{2014}$
在 2014 年及以后取 1，交互项为标准化情绪与该虚拟变量的乘积。

\begin{table}[htbp]
\centering
\small
\begin{threeparttable}
\caption{主要变量定义}
\label{tab:var}
\begin{tabular}{lll}
\toprule
变量 & 定义 & 说明 \\
\midrule
tone\_P        & 政策指引段情绪，标准化                 & 分母为零时记为缺失 \\
tone\_R        & 宏观回顾段情绪，标准化                 & 同上 \\
$R^{(1)}$      & 发布后一日市场收益                     & $I_{t_1}/I_{t_0}-1$ \\
$R^{(5)}$      & 发布后五日市场收益                     & $I_{t_5}/I_{t_0}-1$ \\
CAR            & 发布前后三日累计收益                   & $I_{t_1}/I_{t_{-1}}-1$ \\
$|R^{(0)}|$    & 事件日绝对收益                         & 波动代理之一 \\
Range          & 事件日日内振幅                         & 16 行业 $(H-L)/C_{-1}$ 的均值 \\
PreVol         & 事前 60 日波动                         & 发布前市场收益标准差 \\
$\text{post}_{2014}$ & 2014 年及以后为 1               & 制度断点虚拟变量 \\
\bottomrule
\end{tabular}
\begin{tablenotes}[flushleft]
\footnotesize
\item 收益与振幅均为小数，回归中系数按基点（$10^{-4}$）报告。
\end{tablenotes}
\end{threeparttable}
\end{table}

经上述处理，可用于回归的样本为 99 期，覆盖 2001Q1 至 2026Q2；其中 98 期可计算
日内振幅（1 期因行情缺失）。2014 年及以后的子样本为 50 期。

\paragraph{数据与代码可得性}
报告原文取自人民银行官网公开页面，行业行情取自 akshare 公开接口，均为公开可得数据，
不涉及脱敏问题。情感词典的使用遵循其说明页的引用要求。全部清洗、度量与估计脚本随论文
一并提供，按说明可复现正文表格。
